\documentclass[11pt,a4paper]{article}
\usepackage[utf8]{inputenc}
\usepackage[T1]{fontenc}
\usepackage[italian]{babel}
\usepackage[margin=2.5cm]{geometry}
\usepackage{booktabs}
\usepackage{amsmath}
\usepackage{siunitx}
\usepackage{multirow}
\usepackage{float}
\usepackage[autostyle]{csquotes}
\usepackage[hidelinks]{hyperref}
\usepackage{graphicx}

\sisetup{detect-weight=true}

\title{\bfseries Tesi FL Marocco --- aggiornamento della settimana del 09/11}
\author{}
\date{}

\begin{document}
\maketitle
\vspace{-3.2em}
\begin{center}
\emph{Client-Driven Local Work Assignment in Battery-Constrained Federated Learning}
\end{center}
\vspace{1em}

%======================================================================
\section{Come funziona SAGE-smart 2}
\label{sec:protocol}

\subsection{Round 0: iscrizione (nessun training)}

Ogni telefono si iscrive presso il server e gli invia:
\begin{itemize}
  \item i suoi dati statici: velocità di calcolo (tier), capacità della
        batteria e numero di campioni locali;
  \item il suo stato di carica (SoC) iniziale.
\end{itemize}
Il server
assume per tutti il carico medio della popolazione, $w = 0.118$. Inoltre, dai dati statici il server fissa la \textbf{scadenza del round}: la mediana,
sui telefoni, del tempo necessario per il lavoro locale standard (5 epoche con
$B = 128$). Vale \SI{132}{s} e resta fissa per tutta la run.

\subsection{Ogni round, in sette passi}

\paragraph{1. Previsione (stima del SoC lato server).}
Per ogni telefono il server conserva l'ultimo stato ricevuto e il momento in
cui l'ha ricevuto. Prevede:
\begin{itemize}
  \item la scarica dall'ultimo report;
  \item la probabilità che nel frattempo il proprietario abbia messo il telefono in carica (se era in carica, quanto si è ricaricato).
\end{itemize}

\paragraph{2. Classifica.}
Ogni telefono riceve un punteggio
\begin{equation*}
  \text{punteggio} = a \cdot \widehat{\mathrm{SoC}} + b \cdot D + c \cdot S ,
\end{equation*}

\paragraph{3. Numero di posti.}
Il round $r$ ha $r + 1$ posti: 2 al round 1, 3 al round 2, fino a 20 dal
round 19 in poi. Il tetto è il 20\,\% dei telefoni ancora vivi.

\paragraph{4. Proposte (negoziazione).}
Il server non manda il modello ai primi 20 della classifica: manda una
\textbf{proposta} ai primi $\lceil 1.5 \times \text{posti} \rceil$, cioè
30 telefoni una volta finita la rampa.
\begin{itemize}
  \item Circa il 10\,\% di loro non è raggiungibile (fuori copertura o spento).
  \item Un telefono che risponde invia il suo stato attuale (SoC, carico, se
        è in carica) e accetta o rifiuta.
  \item Un telefono rifiuta solo per energia: quando il lavoro minimo (2 epoche)
        lo porterebbe sotto la riserva del 20\,\% di batteria. Un telefono in
        carica accetta sempre.
  \item Se le accettazioni non riempiono i posti, una seconda ondata di proposte
        va ai telefoni successivi della classifica (al massimo due ondate).
\end{itemize}

\paragraph{5. Scelta fresca.}
Fra i telefoni che hanno accettato, il server tiene i migliori, tanti quanti
i posti, con lo stato che \emph{hanno appena inviato}:
\begin{equation*}
  \text{punteggio fresco} = \bigl(a \cdot \mathrm{SoC}_{\text{ora}} + b \cdot D
                        + c \cdot S\bigr) \cdot \bigl(1 - w_{\text{ora}}\bigr).
\end{equation*}
\begin{itemize}
  \item Un telefono in carica conta come $\mathrm{SoC} = 1$: il training non
        costa nulla alla sua batteria.
  \item Un telefono occupato pesa meno: entro la scadenza farebbe meno passi.
  \item Gli altri telefoni che avevano accettato sono \textbf{rinviati}:
        mantengono la loro anzianità e sono in cima alla classifica nel round
        successivo.
\end{itemize}

\paragraph{6. Envelope e training locale.}
I telefoni scelti ricevono il modello e un \textbf{envelope}: $E$ fra 2 e un
massimo calcolato per quel telefono (al più 10, limitato dalla scadenza con il
suo carico attuale e dalla sua batteria sopra la riserva), $B$ fra 16, 32, 64
e 128, la scadenza (\SI{132}{s}) e la riserva (20\,\%). Poi il telefono:
\begin{itemize}
  \item misura la loss del modello su 256 dei suoi campioni: una loss alta
        significa che i suoi dati sono nuovi per il modello, quindi conviene
        lavorare di più;
  \item sceglie $E$ in proporzione a quella loss e alla batteria che ha sopra
        la riserva;
  \item sceglie $B$ partendo dal più piccolo, perché un batch piccolo dà più
        passi di SGD nello stesso tempo;
  \item se nemmeno 2 epoche stanno nella scadenza, fa solo i passi di SGD che
        ci stanno (lavoro parziale);
  \item rimanda il modello aggiornato e un report: passi fatti, istogramma
        delle classi, SoC finale.
\end{itemize}

\paragraph{7. Aggregazione e aggiornamenti.}
Il server:
\begin{itemize}
  \item fa la media dei modelli con FedAvg, pesando ogni update per i campioni
        visti davvero (un update parziale pesa meno);
  \item aggiorna, per i telefoni che si sono allenati, la divergenza $D$, lo
        stato noto e l'anzianità (che torna a 0);
  \item valuta il modello globale sul test set.
\end{itemize}

\newpage

%======================================================================
\section{Ablazione cumulata di SAGE-smart 2}
\label{sec:ablation}

Partendo da un protocollo ridotto all'osso, a ogni gradino si aggiunge un
componente. Tutte le run usano il seed 2026, la batteria da datasheet, il
carico nominale delle altre app e 250 round; le barre della
Figura~\ref{fig:cum} e i valori della Tabella~\ref{tab:cum} sono la media su
$\beta = 0.1$, $0.5$ e $1.0$.

\begin{itemize}
  \item[0] \textbf{Base}: punteggio $a \cdot \mathrm{SoC} + b \cdot D$, $E = 5$ e
        $B = 128$ fissi, niente negoziazione, SoC dall'ultimo report (nessuna
        stima), FedAvg, broadcast dello stato a tutti i telefoni a ogni round.
  \item[1] + termine di anzianità ($c$).
  \item[2] + $E$ adattivo, scelto dal client dentro l'envelope; comprende il
        lavoro parziale.
  \item[3] + $B$ adattivo, scelto dal client fra 16, 32, 64 e 128.
  \item[4] + negoziazione e stima del SoC lato server insieme, come descritto
        nella Sezione~\ref{sec:protocol}. Questo è SAGE-smart 2.
\end{itemize}

\begin{figure}[H]
  \centering
  \includegraphics[width=\textwidth]{fig_cumulative_ablation_smart2.png}
  \label{fig:cum}
\end{figure}

\begin{table}[H]

  \label{tab:cum}
  \begin{tabular}{lcccc}
    \toprule
    Gradino & Accuratezza (\%) & Dispositivi scarichi & SoC medio finale (\%) & Energia (Wh) \\
    \midrule
    0 base                & 41.9          & 33.7           & 35.7          & 2016 \\
    1 + anzianità         & 47.5 (+5.6)   & 31.0 ($-$2.7)  & 34.5 ($-$1.2) & 2024 (+0.4\,\%) \\
    2 + $E$ adattivo      & 45.0 ($-$2.5) & 3.7 ($-$27.3)  & 38.3 (+3.8)   & 1098 ($-$45.7\,\%) \\
    3 + $B$ adattivo      & 56.5 (+11.5)  & 5.0 (+1.3)     & 39.5 (+1.2)   & 1106 (+0.7\,\%) \\
    4 + negoziazione + stima & 56.1 ($-$0.4) & 3.3 ($-$1.7) & 39.0 ($-$0.5) & 1101 ($-$0.5\,\%) \\
    \midrule
    FedAvg (riferimento)  & 51.1          & 66.0           & 14.4          & 1984 \\
    \bottomrule
  \end{tabular}
\end{table}

\newpage

\subsection{Perché negoziazione e stima del SoC adesso hanno senso}

\begin{table}[H]
  \centering
  \caption{Gradino 3 (broadcast dello stato, niente negoziazione, SoC
  dall'ultimo report) contro gradino 4 (SAGE-smart 2), media sui tre $\beta$.}
  \label{tab:negest}
  \begin{tabular}{lcc}
    \toprule
    & Gradino 3 & Gradino 4 (SAGE-smart 2) \\
    \midrule
    Energia di controllo (Wh per run)                  & 10.3  & 1.0 \\
    Errore del SoC usato per la classifica (punti)     & 1.27  & 0.73 \\
    Telefoni in carica scelti per round                & 1.14  & 1.44 \\
    Proposte rifiutate                                 & ---   & 0.5\,\% \\
    Telefoni rinviati per round (mantengono l'anzianità) & --- & 6.7 \\
    Dispositivi scarichi                               & 5.0   & 3.3 \\
    Accuratezza (\%)                                   & 56.5  & 56.1 \\
    Energia (Wh)                                       & 1106  & 1101 \\
    \bottomrule
  \end{tabular}
\end{table}

\begin{itemize}
  \item \textbf{Niente broadcast.} Il server non contatta più tutti i 100
        telefoni a ogni round solo per sapere chi è raggiungibile: le risposte
        alle proposte danno sia la raggiungibilità sia lo stato. L'energia di
        controllo cala del 90\,\% (da 10.3 a \SI{1.0}{Wh}).
  \item \textbf{Informazione migliore con meno messaggi.} Grazie alla
        previsione e alle risposte, il SoC che il server usa per ordinare i
        telefoni è più vicino a quello vero: 0.73 punti di errore contro 1.27.
  \item \textbf{La scelta finale usa lo stato di adesso.} La scelta fresca evita
        i telefoni occupati in quel momento e preferisce quelli in carica, il cui
        training non costa batteria: si scelgono il 26\,\% di telefoni in carica
        in più (1.44 contro 1.14 per round).
  \item \textbf{Essere interpellato non costa nulla al telefono.} I rifiuti sono
        rari (0.5\,\%), perché un telefono rifiuta solo per energia, e i 6.7
        telefoni rinviati per round mantengono la loro anzianità. È l'opposto di
        SAGE-smart, dove essere interpellato e rifiutare faceva retrocedere un
        telefono.
\end{itemize}
Il risultato è la stessa accuratezza e la stessa energia con il 34\,\% di
dispositivi scarichi in meno, un protocollo più economico e un'informazione più
precisa. La variazione di accuratezza ($-$0.4 punti in media: $-$1.4, $-$0.4 e
+0.6 per i tre $\beta$) è dentro il rumore di un singolo seed: due varianti che
scelgono gli stessi telefoni con lo stesso lavoro nei round 1 e 2 hanno già
un'accuratezza diversa al round 1 (2.02\,\% contro 1.51\,\%), e da lì le due
run seguono traiettorie diverse. Con un solo seed, differenze sotto circa 2
punti fra due varianti non sono significative.

\section{Cose da fare adesso}
\begin{itemize}
    \item
\end{itemize}

\end{document}
